%----------------------------------------------------------------------------------------
%	PACKAGES AND OTHER DOCUMENT CONFIGURATIONS
%----------------------------------------------------------------------------------------

\documentclass[
  12pt,
]{fphw}

\usepackage[utf8]{inputenc}
\usepackage[T1]{fontenc}
\usepackage{mathpazo}
\usepackage{float}
\usepackage{graphicx}
\usepackage{booktabs}
\usepackage{amsmath}
\usepackage{subcaption}
\usepackage{caption}
\usepackage{url}
\usepackage[colorlinks=true, urlcolor=blue]{hyperref}

%----------------------------------------------------------------------------------------
%	ASSIGNMENT INFORMATION
%----------------------------------------------------------------------------------------

\title{Final Project Proposal: RGB-D Object Detection for Collision Avoidance}
\author{Thomas Kulch, Darshan Kedari, Wade Williams}
\date{July 21, 2026}
\institute{Northeastern University // Khoury College of Computer Science}
\class{Computer Vision (CS 5330)}
\professor{Bruce Maxwell}

% Compact title block: smaller title font, tighter spacing, one-line course/date
\renewcommand{\StudentLang}{Students}
\makeatletter
\renewcommand{\maketitle}{%
    \thispagestyle{plain}
    \begin{center}
        \ifx\instituteloc\undefined{}\else{\scshape \instituteloc} \\[4pt] \fi
        \hrule
        \vspace{6pt}
        {\Large \bfseries \@title} \\[3pt]
        \StudentLang: {\itshape \@author}
        \vspace{4pt}
        \hrule
        \vspace{6pt}
        \ifx\classlocal\undefined{}\else\CourseLang: {\itshape \classlocal{} \ifx\professorloc\undefined{}\else{-- }\fi  }\fi
        \ifx\professorloc\undefined{}\else{\ProfessorLang: {\itshape\professorloc\/}}~\fi
        -- \DateLang: {\itshape\@date}
        \vspace{8pt}
    \end{center}
}
\makeatother

%----------------------------------------------------------------------------------------

\begin{document}

\maketitle

%----------------------------------------------------------------------------------------

In our project, we will build a perception pipeline that converts an RGB-D video stream into a top-down
occupancy map with real-world distances and scale for collision avoidance. For each frame, the system will
detect and classify objects in the RGB image, place every obstacle in real-world 3D coordinates using
aligned depth, remove the ground plane so the floor is not treated as an obstacle, and cluster the
remaining 3D points so that obstacles the detector was never trained on are still detected. The output of
this pipeline will be a Bird's Eye View (BEV) occupancy map labeling each obstacle's class and 3D
position. We plan to evaluate the system on at least three datasets --- the KITTI 3D object detection
benchmark~\cite{kitti} for outdoor environments, SUN RGB-D~\cite{sunrgbd} on indoor environments, and on
our own recordings that we will collect using an RGB-D camera. SUN RGB-D and our collected dataset will
come with real aligned depth measured by RGB-D sensors, while KITTI only offers RGB images with LiDAR
measurements as the ground truth. Therefore, to obtain depth data for the KITTI dataset, we will use Depth
Anything V2~\cite{depthanything2} to predict depth, and for the datasets where we already have depth data,
we will run variants using both Depth Anything V2 predictions and the depth data from the RGB-D sensor
data. This one experiment will let us measure to what extent depth estimates from Depth Anything V2 can
substitute real sensor data.

The pipeline will consist of six stages. First, a YOLO11 detector~\cite{yolo}, pretrained on the COCO
dataset~\cite{coco}, will produce 2D boxes and classes, with COCO's classes mapped to each dataset's
labels. It will then back-project the depth map through the camera intrinsics into a 3D point cloud,
serving as a LiDAR-like representation~\cite{pseudolidar}. Next, the pipeline will use
RANSAC~\cite{ransac} to remove the ground-plane as an object and a clustering algorithm to group remaining
points into objects ourselves if needed. We will then fuse the YOLO object detections with the 3D clusters
by projecting the clusters onto the image, so that each object gets a class label (or ``unknown''), a 3D
centroid, and a bounding box. Lastly, we will render obstacles on a BEV occupancy map grid with distance
rings / shapes.

We will evaluate the success of our pipeline in various ways depending on the dataset. On KITTI, we will
compute the the 3D centroid localization error against the dataset's 3D box annotations to measure how
accurately we place objects and 3D average precision from the KITTI evaluation code to measure how
reliably we find them. We'll compare both metrics to published pseudo-LiDAR results~\cite{pseudolidar}. On
SUN RGB-D, we will evaluate the same two metrics for the dataset's sensor depth -- mean average precision
under the community-standard 10-class protocol~\cite{votenet} and the centroid error. The COCO dataset
only covers part of this 10-class list, so we will evaluate the classes that are scoreable, with the rest
appearing as unknown obstacles. We may also interchange the nuScenes dataset~\cite{nuscenes} for the KITTI
dataset as outdoor benchmarks depending on our preliminary results.

Beyond this core system, we have brainstormed possible extensions we could pursue depending on our initial
progress and results by the checkin date. These extensions include recording our own data using an RGB-D
camera with a UI that shows detected objects and collision alerts and recording a live demo, utilizing our
pipeline to enable collision steering via object detection in the Habitat-Sim simulation
environment~\cite{habitat} where a mobile robot navigates photorealistic indoor environments using the BEV
map, finetuning the object detector on target domains, finetuning the metric depth head on KITTI's ground
truth LiDAR depth measurements, and training a frustum-based 3D regression head~\cite{frustum}, with the
results reported as an ablation against the zero-shot baseline. Further, we will consider implementing our
pipeline on a real robot and experimenting with 3D scene reconstruction methods, but this is likely out of
scope. We are very excited about this project and look forward to hearing feedback and discussing
preliminary results.

%----------------------------------------------------------------------------------------

\begin{thebibliography}{11}

\bibitem{kitti}
A.~Geiger, P.~Lenz, and R.~Urtasun, ``Are we ready for autonomous driving? The KITTI vision benchmark
suite,'' in \emph{CVPR}, 2012.

\bibitem{sunrgbd}
S.~Song, S.~Lichtenberg, and J.~Xiao, ``SUN RGB-D: A RGB-D scene understanding benchmark suite,'' in
\emph{CVPR}, 2015.

\bibitem{depthanything2}
L.~Yang, B.~Kang, Z.~Huang, Z.~Zhao, X.~Xu, J.~Feng, and H.~Zhao, ``Depth Anything V2,'' in
\emph{NeurIPS}, 2024.

\bibitem{yolo}
J.~Redmon, S.~Divvala, R.~Girshick, and A.~Farhadi, ``You only look once: Unified, real-time object
detection,'' in \emph{CVPR}, 2016.

\bibitem{coco}
T.-Y.~Lin, M.~Maire, S.~Belongie, J.~Hays, P.~Perona, D.~Ramanan, P.~Doll\'ar, and C.~L.~Zitnick,
``Microsoft COCO: Common objects in context,'' in \emph{ECCV}, 2014.

\bibitem{pseudolidar}
Y.~Wang, W.-L.~Chao, D.~Garg, B.~Hariharan, M.~Campbell, and K.~Q.~Weinberger, ``Pseudo-LiDAR from visual
depth estimation: Bridging the gap in 3D object detection for autonomous driving,'' in \emph{CVPR}, 2019.

\bibitem{ransac}
M.~A.~Fischler and R.~C.~Bolles, ``Random sample consensus: A paradigm for model fitting with applications
to image analysis and automated cartography,'' \emph{Communications of the ACM}, vol.~24, no.~6, 1981.

\bibitem{votenet}
C.~R.~Qi, O.~Litany, K.~He, and L.~J.~Guibas, ``Deep Hough voting for 3D object detection in point
clouds,'' in \emph{ICCV}, 2019.

\bibitem{nuscenes}
H.~Caesar, V.~Bankiti, A.~H.~Lang, S.~Vora, V.~E.~Liong, Q.~Xu, A.~Krishnan, Y.~Pan, G.~Baldan, and
O.~Beijbom, ``nuScenes: A multimodal dataset for autonomous driving,'' in \emph{CVPR}, 2020.

\bibitem{habitat}
M.~Savva, A.~Kadian, O.~Maksymets, Y.~Zhao, E.~Wijmans, B.~Jain, J.~Straub, J.~Liu, V.~Koltun,
J.~Malik, D.~Parikh, and D.~Batra, ``Habitat: A platform for embodied AI research,'' in \emph{ICCV}, 2019.

\bibitem{frustum}
C.~R.~Qi, W.~Liu, C.~Wu, H.~Su, and L.~J.~Guibas, ``Frustum PointNets for 3D object detection from RGB-D
data,'' in \emph{CVPR}, 2018.

\end{thebibliography}

\end{document}
